\section{Results}
\label{sec:results}

% ZN notes
% - draft result figures cover the key bits, write to them
% - comparison across scenarios and vs. CMIP6 probably most interesting bit
% - plot IGCC too where available (shows which scenario is closest to 2023, 2024 and 2025)
% - write then get feedback, hard to guess what people find interesting
%
% - change name of groups so they don't match scenario names
%
% - CO2
%     - roughly even spread and similar shapes to 2100: Either peak and decline or ongoing rise (H, M)
%     - post 2100 settle around 300-400ppm or slow stabilisation consistent with approx. net zero emissions (need to check underlying emissions timeseries and probably put in appendix)
%     - make no comment on separation, that's a different thing determining when signal comes out of noise
%     - compared to CMIP6, top end much lower
%     - compared to CMIP6, medium and low similar range
%     - compared to CMIP6, LN completely new shape and base (well below pre-industrial CO2, due to neg. emms c/f emissions)
%     - VL-CF lowering peak CO2 by Xppm
%     - LN who knows what happens to lat. gradient. Great one to look at once CMIP7 esm- results come out
%
% - CH4
%     - low group consistent with decline in methane emissions, just different floor
%     - H, M and HL all similar rise, then question of when rise stops. HL then plummets, others roughly stabilise
%     - post 2100 settle around 1000ppb or stabilisation consistent with approx. constant emissions (need to check underlying emissions timeseries and probably put in appendix)
%     - make no comment on separation, that's a different thing determining when signal comes out of noise
%     - compared to CMIP6, top end slightly lower although HL is a new high peak
%     - compared to CMIP6, medium similar but different extension
%     - compared to CMIP6, low has higher peak, better spread and different extension (check emissions, CMIP6 had emissions that went back to zero I believe)
%     - VL-CF CH4 similar to VL. Makes sense as emissions are the same. You're seeing feedback (and harmonisation up until that has converged) effects here
%
% - N2O
%     - roughly even spread
%     - gap between VL and the rest (VL works hard on N2O)
%     - post 2100 settle all decline back towards pre-industrial (need to check underlying emissions timeseries and probably put in appendix)
%     - VL-CF makes harmonisation effect clear: MAGICC hasn't cpatured 'kick' since 2015 and this is obvious in VL-CF. N2O is trickiest for harmonisation, most important next fix
%     - make no comment on separation, that's a different thing determining when signal comes out of noise
%     - compared to CMIP6, very similar
%     - compared to CMIP6, VL is new shape at low end
%     - compared to CMIP6, extensions completely different. CMIP6 used constant ag. emissions. CMIP7 does not. No obvious answers of correct but clear implications of choice (will need plots to show this)
%     - VL-CF CH4 similar to VL. Makes sense as emissions are the same. You're seeing feedback and harmonisation up until that has converged effects here, which demonstrates how big the harmonisation issue can be when model doesn't match history well.
%
% - N2O
%     - roughly even spread
%     - gap between VL and the rest (VL works hard on N2O)
%     - post 2100 settle all decline back towards pre-industrial (need to check underlying emissions timeseries and probably put in appendix)
%     - VL-CF makes harmonisation effect clear: MAGICC hasn't cpatured 'kick' since 2015 and this is obvious in VL-CF. N2O is trickiest for harmonisation, most important next fix
%     - make no comment on separation, that's a different thing determining when signal comes out of noise
%     - compared to CMIP6, very similar
%     - compared to CMIP6, VL is new shape at low end
%     - compared to CMIP6, extensions completely different. CMIP6 used constant ag. emissions. CMIP7 does not. No obvious answers of correct but clear implications of choice (will need plots to show this)
%     - VL-CF CH4 similar to VL. Makes sense as emissions are the same. You're seeing feedback and harmonisation up until that has converged effects here, which demonstrates how big the harmonisation issue can be
%
% - CFC12-eq
%     - ongoing decline projected
%     - slightly lower than CMIP6 now, slightly higher long-term (due to better handling of slow release I think?)
%
% - HFC134a-eqq
%     - all over the place, small effect (quantify ERF) so usually not modelled carefully (plot emissions too). ML and M shape 'creative'. L probably out of line with narrative
%     - big spread. H very high
%     - extensions consistent with decline in emissions (check plots), different to CMIP6 which went with constant I think (CMIP7 better here probably)
%     - slightly lower than CMIP6 now, slightly higher long-term (due to better handling of slow release I think?)
%     - compared to CMIP6, lower top end
%     - VL-CF matches VL exactly
